\section{Give the full-byte tags priority}
At a command boundary, read one tag byte after proving it is inside the payload. Compare
the complete byte with \texttt{fe} and \texttt{ff} first. They select RGB and RGBA
commands. Only after excluding them should the decoder examine the high two bits.
Those bits select INDEX, DIFF, LUMA or RUN for the remaining bytes.

\tagpdfsetup{table/header-rows={1}}
\begin{center}
  \begin{tabular}{lll}
    Tag                                         & Additional bytes & Reconstruction                \\
    \hline
    \texttt{fe}                                 & Three            & Replace RGB, retain alpha     \\
    \texttt{ff}                                 & Four             & Replace RGBA                  \\
    \texttt{00xxxxxx}                           & None             & Fetch one cache entry         \\
    \texttt{01xxxxxx}                           & None             & Apply three small differences \\
    \texttt{10xxxxxx}                           & One              & Apply green and two residuals \\
    \texttt{11xxxxxx}, excluding \texttt{fe,ff} & None             & Repeat one through 62 pixels  \\
  \end{tabular}
\end{center}

Masking \texttt{fe} with \texttt{c0} would classify it as a run. That error invents
sixty-three output pixels and leaves its three actual samples to be parsed as tags.
The resulting fault is larger than one wrong color: it destroys both cursors' meanings.
A full-byte comparison before the two-bit dispatch prevents the ambiguity.

For RGB, prove that three payload bytes remain, then replace red, green and blue while
retaining the previous alpha. For RGBA, require four and replace every component.
Neither command performs a colorspace transform or premultiplies RGB by alpha.
